\documentclass[12pt,letterpaper]{article}

\usepackage{iftex}
\ifPDFTeX
  \usepackage[T2A]{fontenc}
  \usepackage[utf8]{inputenc}
\else
  \usepackage{fontspec}
  \setmainfont{cmunrm.otf}[
    BoldFont=cmunbx.otf,
    ItalicFont=cmunti.otf,
    BoldItalicFont=cmunbi.otf]
  \setmonofont{cmuntt.otf}
\fi
\usepackage[english,russian]{babel}
\usepackage{amsmath,amssymb}
\usepackage[margin=1in]{geometry}
\usepackage{enumitem}
\usepackage{fancyhdr}
\usepackage{microtype}

\setlength{\parindent}{0pt}
\setlength{\parskip}{0pt}
\setlength{\headheight}{15pt}
\setlist[enumerate]{label=(\alph*),leftmargin=3em,labelsep=0.5em,
  topsep=8pt,itemsep=7pt,parsep=0pt}
\pagestyle{fancy}
\fancyhf{}
\fancyhead[R]{\thepage}
\renewcommand{\headrulewidth}{0.4pt}
\fancypagestyle{plain}{
  \fancyhf{}
  \fancyfoot[C]{\thepage}
  \renewcommand{\headrulewidth}{0pt}
}

\newcounter{task}
\newcommand{\task}{\par\addvspace{0.8\baselineskip}%
  \refstepcounter{task}\noindent\textbf{Задача \thetask.}\enspace}
\DeclareMathOperator{\softmax}{softmax}
\DeclareMathOperator{\Att}{Att}
\DeclareMathOperator{\stopgrad}{stopgrad}

\title{Глубинное обучение 1, ФКН ВШЭ\\
Теоретическое домашнее задание №3\\
Эмбеддинги и языковые модели}
\author{}
\date{}

\begin{document}
\maketitle

\task
Пусть для корпуса слов \texttt{a b a c} обучается модель Word2Vec
skip-gram с радиусом контекста $K=1$. Каждому слову $w$ соответствуют
\emph{два} вектора: центральный $u_w\in\mathbb{R}^2$ и контекстный
$v_w\in\mathbb{R}^2$. В словаре дополнительно есть слово \texttt{d},
отсутствующее в данном предложении.
\begin{enumerate}
  \item Выпишите все пары «центральное слово~--- слово контекста»,
  получаемые из корпуса, \emph{с учётом кратностей}.

  \item Для положительной пары $(a,b)$ возьмём один отрицательный
  пример $d$. Функция потерь для этой пары:
  \[
    \ell=-\ln\sigma(u_a^\top v_b)-\ln\sigma(-u_a^\top v_d),
    \qquad \sigma(x)=\frac{1}{1+e^{-x}}.
  \]
  Пусть $u_a=(1,0)^\top$, $v_b=(0,1)^\top$, $v_d=(1,0)^\top$.
  Найдите $\ell$, $\nabla_{u_a}\ell$ и новый вектор $u_a$ после одного
  шага градиентного спуска с $\eta=\tfrac12$ (векторы $v_b,v_d$
  на этом шаге не изменяются). Проверьте, как изменились скалярные
  произведения $u_a^\top v_b$ и $u_a^\top v_d$.

  \item Докажите, что при любом обратимом $A\in\mathbb{R}^{2\times2}$
  одновременная замена
  \[
    u_w\longmapsto Au_w,\qquad v_w\longmapsto A^{-\top}v_w
  \]
  не меняет ни одно скалярное произведение $u_w^\top v_{w'}$ и,
  следовательно, не меняет функцию потерь skip-gram с negative sampling.
  Обязательно ли при этом сохраняются евклидовы расстояния между
  центральными эмбеддингами? Докажите ответ или приведите контрпример.
\end{enumerate}

\task
Рассмотрим скалярную RNN без нелинейности. При дифференцировании
считайте $h_0$ переменной, а производные вычисляйте в точке $h_0=1$:
\[
  h_0=1,\qquad h_t=ah_{t-1}\quad(t=1,2,3,4),\qquad L=\tfrac12 h_4^2.
\]
\begin{enumerate}
  \item Выразите $L$ через $a$, найдите точные производные
  $\frac{\partial L}{\partial a}$ и $\frac{\partial L}{\partial h_0}$.
  Вычислите их при $a=\tfrac12$ и $a=2$ и объясните разницу в масштабах
  градиентов.

  \newpage

  \item Теперь при обратном проходе \emph{оборвём граф} между шагами
  $t=2$ и $t=3$: в вычислениях $h_3$ используется $\stopgrad(h_2)$,
  численно равное $h_2$, но не пропускающее через себя градиент.
  Найдите градиенты $\frac{\partial L}{\partial a}$ и
  $\frac{\partial L}{\partial h_0}$ в этом случае. Сравните их
  с пунктом a) и объясните, почему общий параметр $a$ всё же получает
  ненулевой градиент.

  \item В упрощённой рекуррентной ячейке
  $h_t=(1-z)h_{t-1}+zc_t$ число $z\in(0,1)$ постоянно,
  а $c_t$~--- константы. Найдите $\frac{\partial h_T}{\partial h_0}$
  и сравните его при $T=4$, $z=0.1$ с производной
  $\frac{\partial h_4}{\partial h_0}$ исходной RNN при $a=\tfrac12$.
  Что в этом примере меняет наличие гейта?
\end{enumerate}

\task
Для одной головы внимания с $d_k=1$ заданы уже вычисленные проекции
\[
  Q=\begin{pmatrix}1\\0\\1\end{pmatrix},\qquad
  K=\begin{pmatrix}0\\\ln 2\\\ln 2\end{pmatrix},\qquad
  V=\begin{pmatrix}1\\3\\5\end{pmatrix}.
\]
Выход слоя $Z=\softmax(QK^\top+M)V$, где $\softmax$ применяется
к каждой строке. Масштабирующий коэффициент $1/\sqrt{d_k}=1$.
Для \emph{немаскированного} внимания $M=0$; для \emph{причинного}
внимания $M_{ij}=0$ при $j\leq i$ и $M_{ij}=-\infty$ при $j>i$
(индексация с единицы).
\begin{enumerate}
  \item Для обоих вариантов вычислите матрицу $S$ и вектор $Z$.
  Здесь $S=\softmax(QK^\top+M)$. Ответы дайте в виде точных дробей.

  \item Считая $Q$ и $K$ фиксированными, найдите
  $\frac{\partial Z_1}{\partial V_3}$ в обоих вариантах. Объясните,
  почему результаты различаются и как это связано с авторегрессионной
  генерацией.

  \item Пусть $P$~--- любая перестановочная матрица размера $3\times3$.
  Докажите, что \emph{без маски и позиционных кодировок}
  \[
    \Att(PQ,PK,PV)=P\Att(Q,K,V).
  \]
  Верно ли это равенство с фиксированной причинной маской для всех $P$?
  Обоснуйте ответ.
\end{enumerate}

\task
Рассмотрим дистилляцию для двух классов. Исходные логиты учителя
и студента: $z^T=(\ln 4,0)$, $z^S=(0,0)$. Для температуры $\tau>0$
определим
\[
  p_\tau=\softmax(z/\tau),\qquad
  \ell_\tau=-p^T_{\tau,1}\ln p^S_{\tau,1}
            -p^T_{\tau,2}\ln p^S_{\tau,2}.
\]
Учитель фиксирован. Производные берутся только по логитам студента;
температура и распределение учителя при дифференцировании не меняются.
\begin{enumerate}
  \item Вычислите распределения учителя и студента при $\tau=1$
  и $\tau=2$. Найдите $\ell_1$ и $\ell_2$ в исходной точке.
  Ответы дайте в точном виде.

  \newpage

  \item Выведите градиент $\ell_\tau$ по $z^S$. Вычислите его при
  $\tau=1$ и $\tau=2$ в исходной точке. Для каждой температуры независимо
  сделайте шаг градиентного спуска с $\eta=1$ по логитам студента.
  Объясните, почему одинаковые значения потерь не означают одинаковые
  градиенты.

  \item Пусть истинная метка~--- класс 2, $\tau=1$, а общая потеря равна
  $L_a=(1-a)(-\ln p^S_{1,2})+a\ell_1$, $a\in(0,1)$.
  Найдите оптимальное распределение студента при $a=1/2$ и определите,
  при каких $a$ он предсказывает класс 1. Случай равных вероятностей
  разберите отдельно.
\end{enumerate}

\task
Датасет состоит из последовательностей длины два:
\[
  (X_1,X_2)=(Z,Z),\qquad
  \mathbb{P}(Z=0)=\mathbb{P}(Z=1)=\tfrac12.
\]
Рассмотрим \emph{упрощённый} masked language modeling: каждая из двух
позиций независимо выбирается с вероятностью $p\in(0,1)$; выбранный
токен всегда заменяется на \texttt{[MASK]} (правило 80/10/10 в этой
задаче \emph{не используется}). Функция потерь~--- сумма
$-\ln P_\theta(X_i\mid\text{искажённый вход})$ \emph{только по выбранным
позициям}. Предположим, модель имеет достаточно параметров, чтобы
реализовать оптимальные условные вероятности.
\begin{enumerate}
  \item Рассмотрите четыре возможных множества выбранных позиций:
  $\varnothing$, $\{1\}$, $\{2\}$, $\{1,2\}$. Для каждого укажите,
  что наблюдает двунаправленный энкодер, и найдите минимально достижимую
  сумму потерь на примере при этом варианте маскирования.

  \item Найдите математическое ожидание минимальной MLM-потери для
  \emph{двунаправленного} энкодера через $p$ и вычислите его при
  $p=\tfrac12$. Может ли это ожидание равняться нулю при $p>0$? Почему?

  \item Теперь вместо двунаправленного внимания используется
  \emph{причинное}: при предсказании позиции $i$ разрешено видеть только
  искажённые позиции $1,\ldots,i$. Найдите минимально возможное
  математическое ожидание той же MLM-потери. Сравните с пунктом b)
  и объясните, откуда возникает разница.
\end{enumerate}

\end{document}
